\chapter{Filtering Empty Camera-Trap Images in Embedded
Systems}\label{cap:filtering}

In this chapter, we present a comparative study on empty frame detection in
camera-trap images to analyze the trade-off between precision and inference
latency on edge devices. To accomplish this objective, we investigate
classifiers and object detectors of various input resolutions and optimize them
using quantization and reducing the number of model filters.

\section{Introduction}

Monitoring wildlife through camera traps produces a massive amount of images,
whose a significant portion does not contain animals, being later discarded.
For instance, about 75\% of the images from the Snapshot Serengeti
dataset do not contain animals~\citep{dryad_5pt92}.

Several studies have investigated the use of deep learning
models to recognize animals in camera trap images~\citep{beery2018recognition,
norouzzadeh2018automatically, schneider2020three, tabak2019machine,
villa2017towards, willi2019identifying}. Due to the development of new
technologies for this type of monitoring,embedding models to identify animals
directly in the devices may provide advantages. For instance, in projects whose
equipment is connected to a network, prior filtering avoids the unnecessary
transmission of empty images, saving network bandwidth and
energy~\citep{elias2017s}. Another example are traditional approaches, which
usually involve going to the capture points to retrieve cameras and/or memory
cards~\citep{team_protocol}. In this case, discarding empty images can save data
storage, extending the time a camera can stay in the field collecting images
without needing maintenance.

Although there are studies comparing the performance of several modern deep
learning architectures on recognizing animals in camera trap
images~\citep{norouzzadeh2018automatically, schneider2020three,
villa2017towards}, there are few studies regarding the trade-off between model
performance and inference latency directly on edge
devices~\citep{tyden2020edge,zualkernan2020towards,schneider2020three}. Thus,
this chapter presents a comparative study between classifiers and detectors of
different complexities running on edge devices to recognize nonempty images.

Our main findings are as follows:
\begin{itemize}
 \item Detection models generally outperform classifiers with
comparable latency if they are trained on the same camera trap image set.
 \item Classifiers can obtain high performance on empty image identification,
 given the difficulty of generating labels to train
detectors. The performance, however, depends on the amount of images
available for training, as well as on specific factors that can vary from
dataset to dataset.
 \item Models with a higher input resolution perform better. Their inference
latency can be kept low by reducing the number of model filters at an
acceptable
performance cost for the problem.
\end{itemize}

\section{Materials and Methods}
Two datasets are investigated in this chapter: 1) Caltech Camera Traps; and 2)
Snapshot Serengeti (SS). These datasets are described below.

\subsection{Datasets}

\textbf{Caltech Camera Traps~\citep{beery2018recognition}:} It contains 243,100
images taken from 140 capture locations in the Southwestern United States. The
instances were labeled in 22 categories, and about 66,000 bounding boxes
localizing animals were also provided. For our experiments, we selected a
subset of images from the empty and nonempty classes, grouping into the nonempty
class images from all other categories that have bounding boxes. Then,
the dataset was partitioned into training and validation set according to the
locations recommended in~\citep{caltech_lila}. Moreover, a subset of the
training partition was split, consisting of 20 locations chosen at random, to
be used to adjust the training hyper-parameters (called here validation dev).
Due to the fact that some locations concentrate a large number of images of the
empty class, which may lead the models to be biased in certain backgrounds, we
decided to limit to 1,000 instances per location the number of empty class
instances in both training and validation dev partitions.
Table~\ref{table:caltech_instances} summarizes the number of instances obtained.

\begin{table}[htb]
\caption[Number of images used from the Caltech dataset]{Number of images used
from the Caltech dataset. The empty class of training and validation dev
partitions was limited to 1000 instances per location. The nonempty class is
composed of images with bounding boxes from all other categories.}
\begin{center}
\begin{tabular}{l|c|c|c}
\hline
\textbf{Class}   & \textbf{Training} & \textbf{Val-dev} & \textbf{Validation}
\\
\hline \hline
Empty    & 8,574            & 2,824              & 19,892              \\
Nonempty & 32,032           & 3,877              & 23,410              \\
\hline
\end{tabular}
\end{center}
\label{table:caltech_instances}
\end{table}

\textbf{Snapshot Serengeti~\citep{dryad_5pt92}:} This dataset currently contains
more than 7 million camera trap images collected over 11 seasons in the
Serengeti National Park, Tanzania. In this work, we use images from the first
six seasons in order to keep consistency with previous works using the
Snapshot Serengeti dataset~\citep{norouzzadeh2018automatically,
villa2017towards, willi2019identifying}. These images were divided into training
and validation sets according to the locations (SS-Site), as recommended
in~\citep{ss_lila}. As was done for the Caltech dataset, we have also created  a
validation dev split for hyper-parameter adjustment. This partition was obtained
by randomly selecting 23 locations from the training set. In addition to the
SS-Site configuration, we have also performed a partitioning by time (SS-Time)
taking into account that the models could be used in images from new seasons. To
generate SS-Time, the first four seasons were grouped into a training set,
while the fifth season was used as validation dev, and the sixth season as the
validation set. Finally, the dataset was adapted to represent a  bi-class
problem, where instances from the blank category were labeled to the empty
class and the others were used to compose the nonempty class. We also balanced
the classes for the training partition.

In addition, since there are approximately only 78,000 images from the nonempty
class annotated with bounding boxes, we also perform experiments using subsets
of instances (called small), which are balanced for the training and validation
dev partitions. The objective here is to perform a fair comparison between
classifiers and detectors when both are trained using the same number of images.
Unless otherwise specified, the experiments conducted in this work use these
subsets with fewer images. Details about all data partitions of the Snapshot
Serengeti dataset investigated in this work are shown in
Table~\ref{table:serengeti_instances}.

\begin{table}[htb]
\caption{Number of instances used from the Snapshot Serengeti dataset.}
\begin{center}
\setlength{\tabcolsep}{4pt}
\begin{tabular}{ll|c|c|c}
\hline
                                                                            &
\textbf{Class}   & \multicolumn{1}{l|}{\textbf{Training}} &
\multicolumn{1}{l|}{\textbf{Val(dev)}} &
\multicolumn{1}{l}{\textbf{Validation}}
\\ \hline \hline
\multirow{2}{*}{SS-Site}                                                    &
Empty    & 524,804                               & 278,578

                 & 535,817                                 \\
                                                                            &
Nonempty & 523,891                               & 84,531

                 & 209,183                                 \\ \hline
\multirow{2}{*}{\begin{tabular}[c]{@{}l@{}}SS-Site\\
(small)\end{tabular}} &
Empty             & 51,281                                & 6,041

                 & 535,817                                 \\
                                                                            &
Nonempty          & 52,081                                & 6,041

                 & 209,183                                 \\ \hline \hline
\multirow{2}{*}{SS-Time}                                                    &
Empty             & 516,635                               & 588,406

                 & 383,981                                 \\
                                                                            &
Nonempty          & 516,630                               & 225,647

                 & 75,328                                  \\ \hline
\multirow{2}{*}{\begin{tabular}[c]{@{}l@{}}SS-Time\\
(small)\end{tabular}} &
Empty             & 43,612                                & 18,957

                 & 383,981                                 \\
                                                                            &
Nonempty          & 44,579                                & 18,957

                 & 75,328                                  \\ \hline
\end{tabular}
\end{center}
\label{table:serengeti_instances}
\end{table}

\subsection{Architectures}

The MobileNetV2~\citep{sandler2018mobilenetv2} architecture is a natural choice
to be used as a classifier baseline, since it was designed specifically for
computationally limited devices. Besides the original MobileNetV2 model (input
resolution $224 \times 224$), a version with input $320 \times 320$ was also
used in our experiments due to the fact that animals can appear very far from
the camera and visually small as a
consequence~\citep{norouzzadeh2018automatically}. From the family of models EfficientNet~\citep{tan2019efficientnet}, we chose two
versions with input resolutions comparable to MobileNetV2's: EfficientNet-B0
($224 \times 224$) and EfficientNet-B3 ($300 \times 300$).

For the detectors, we investigate the SSDLite~\citep{sandler2018mobilenetv2}
with a MobileNetV2 ($320 \times 320$ input) as backbone. For the family of
detectors using EfficientNets as backbones~\citep{tan2020efficientdet}, we
included in our experiments the EfficientDet-D0, which works with an input
resolution of $512 \times 512$, the smallest of this group of models.

\subsection{Implementation Details}

\textbf{Image preprocessing for classification:} We chose a standard procedure
for image preprocessing. Initially, a random rectangular crop of the image is
applied with aspect ratio and area sampled in $[3/4, 4/3]$ and $[65\%, 100\%]$,
respectively. Then, each image is scaled to the input size of each architecture
and a horizontal flip is applied with $50\%$ probability. Next, we apply data
augmentation using RandAugment~\citep{cubuk2020randaugment} with parameters
$N=2$ and $M=2$. Finally, the pixel values of the image are scaled in $[-1, 1]$
for MobileNetV2 and in $[0, 1]$ for EfficientNets. During validation, the
preprocessing consists of only image resizing and pixel scaling depending on
the architecture.

\textbf{Classifiers training procedure:} Each model was initialized with
ImageNet pre-trained weights and then trained during 10 epochs using the
Stochastic gradient descent (SGD) on an NVIDIA GeForce GTX 1080 Ti graphics
card. The initial learning rate was 0.01 for a batch size of 256 and scaled
linearly according to the batch size effectively used, which varied according
to the model due to the graphics card limited memory, as shown in
Table~\ref{table:classifiers_hparams}. The learning rate was linearly increased
from 0 to the initial rate during 30\% of the steps of the first epoch and then
reduced using the cosine decay scheduling, as suggested by \citet{he2019bag}.

\begin{table}[htb]
\caption[Batch size and initial learning rate used for each classifier.]{Batch
size and initial learning rate used for each classifier. The initial learning
rate is scaled linearly according to the batch size $b$, defined by $0.01 \times
b /256$.}
\begin{center}
\begin{tabular}{l|c|l}
\hline
\multicolumn{1}{c|}{\textbf{Architecture}} &
\textbf{\begin{tabular}[c]{@{}c@{}}Batch size\end{tabular}} &
\textbf{\begin{tabular}[c]{@{}c@{}}Learning rate\end{tabular}} \\
\hline \hline
MobileNetV2 (224)                         & 128

                               & 0.005

                       \\
MobileNetV2 (320)                         & 64

                               & 0.0025

                       \\
EfficientNet-B0                           & 32

                               & 0.00125

                       \\
EfficientNet-B3                           & 16

                               & 0.000625

                       \\ \hline
\end{tabular}
\end{center}
\label{table:classifiers_hparams}
\end{table}

\textbf{Detectors training procedure:} Both detectors were initialized with
weights pre-trained on COCO~\citep{lin2014microsoft} and then trained on each
dataset using the Tensorflow Object Detection API~\citep{huang2017speed}. We
have used the standard training procedure for each detector, besides $32$ and
$8$ as batch size for \mbox{SSDLite+MobileNetV2} and EfficientDet-D0
respectively. The learning rate and the number of training epochs were adjusted
according to the performance measured on the validation dev partition. These
values are shown in Table~\ref{table:detectors_hparams}.

\begin{table}[htb]
\caption[Training hyperparameters for detectors.]{Training hyperparameters for
detectors. The same hyperparameters were used for both time and site
partitioning of Snapshot Serengeti (SS).}
\begin{center}
\begin{tabular}{l|l|c|c}
\hline
\textbf{Model} & \textbf{Dataset} &
\textbf{\begin{tabular}[c]{@{}c@{}}Learning\\rate\end{tabular}} &
\textbf{\begin{tabular}[c]{@{}c@{}}Training\\steps\end{tabular}} \\
\hline \hline
\multirow{2}{*}{SSDLite+MobileNetV2}  & Caltech & 0.008

                          & 12,000

           \\
                                 & SS      & 0.004

                          & 36,000

           \\ \hline
\multirow{2}{*}{EfficientDet-D0} & Caltech & 0.008

                          & 50,000
           \\
                                 & SS      & 0.001

                          & 150,000
           \\ \hline
\end{tabular}
\end{center}
\label{table:detectors_hparams}
\end{table}

\subsection{Evaluation Procedure}

In order to compare detectors and classifiers, predictions related to the
bounding box coordinates were ignored. Thus, only the detection confidence
value that represents the class with the highest score was used as the detected
label.

Taking into account that the confidence threshold can be adjusted to avoid
nonempty images discarding, we decided to use the precision-recall curve as
a graphical tool to compare the general performance of the models considering
all possible thresholds. To compare the models effectiveness when discarding
empty images, i.e., the true negative rate (TNR), the confidence threshold
of each model was adjusted so that the recall for nonempty images was set to
96\% -- a value reached by models investigated
in~\citep{norouzzadeh2018automatically}.

We used a Raspberry Pi 3 model B running Raspbian GNU/Linux 10 as a reference
edge device to assess the models' latency. The models were converted to the
TensorFlow Lite format without quantization and their latency was calculated
using the native benchmark tool provided by
TensorFlow\footnote{\url{https://www.tensorflow.org/lite/performance/measurement
}}. The reported latency values were calculated as the average over 50 runs for
each model. Additionally, we also evaluate the models version obtained as a
result of post-training quantization~\citep{jacob2018quantization}. In this
scenario, a subset of 500 training instances was employed to calibrate the
operations to work with the integer type, while maintaining model inputs and
outputs as floating point to keep the original model interface.

\section{Experimental Results}

\subsection{Classifiers vs. Detectors}

In order to compare the performance of classifiers and detectors as fairly as
possible, the models were trained using subsets composed by instances of the
nonempty class annotated with bounding boxes. Despite of the fact that these
subsets were generated using a much smaller amount of the total images
available for classification, leading to possible sub-optimal models, this
number of instances can provide a more realistic perspective of the problem.
Indeed, as pointed out by \citet{schneider2020three}, the vast majority of
small-scale research projects focused on camera trap do not have a large amount
of labeled images.

\textbf{Results:} Figure~\ref{fig:pr_curve} shows the precision-recall curves of
the investigated models for each dataset. As expected, detectors outperformed
classifiers in all datasets, especially Caltech. Considering a confidence
threshold producing 96\% of recall, EfficientDet-D0 was able to eliminate more
than twice as many empty images when compared to the classifiers using Caltech
dataset. In terms of SSDLite+MobileNetV2, it also obtained a significantly
higher true negative rate (at least 19\% in absolute values), as reported in
Table~\ref{table:precision_tvn}. The scenario is quite similar for the Snapshot
Serengeti dataset, since detectors also outperformed classifiers by a
significant margin (at least 8\% of precision). Figure~\ref{fig:predictions}
shows some images to illustrate the results attained in our experiments.


\begin{table}[htb]
\caption[Comparison of precision for the nonempty class and the true negative
rate (TNR)]{Comparison of precision for the nonempty class and the true negative
rate (TNR) where the confidence threshold of each model was adjusted
to achieve a recall of 96\% on the nonempty class. The reported CPU latency
corresponds to the Caltech dataset, but it is similar to the others, being
calculated from the average of 50 runs. The true negative rate indicates the
percentage of images without animals that would no longer be stored
unnecessarily.}
\begin{center}
\setlength{\tabcolsep}{3pt}
\resizebox{\textwidth}{!}{
\begin{tabular}{l|c|ccc|ccc|ccc}
\hline
\textbf{}                            &                &
\multicolumn{3}{c|}{\textbf{Caltech}}                  &
\multicolumn{3}{c|}{\textbf{SS-Site}}                  &
\multicolumn{3}{c}{\textbf{SS-Time}}                  \\
\multicolumn{1}{c|}{\textbf{Model}} & \textbf{CPU}   & \textbf{Precision} &
\textbf{TNR}     & \textbf{Thresh.} & \textbf{Precision} & \textbf{TNR}     &
\textbf{Thresh.} & \textbf{Precision} & \textbf{TNR}   & \textbf{Thresh.}\\
\hline \hline
Efficientnet-B0                      & 801ms          & \textbf{60.26\%}  &
\textbf{25.50\%} & 0.355           & 50.32\%           & 63.00\%          &
0.153           & 33.08\%           & 61.90\%          & 0.159           \\
Efficientnet-B3                      & 3205ms         & 56.42\%           &
12.75\%          & 0.305           & 57.21\%           & 71.97\%          &
0.155           & \textbf{39.27\%}  & \textbf{70.87\%} & 0.170           \\
MobileNetV2-224                      & \textbf{324ms} & 58.60\%           &
20.18\%          & 0.228           & 57.72\%           & 72.55\%          &
0.188           & 30.51\%           & 57.11\%          & 0.126           \\
MobileNetV2-320                      & 638ms          & 58.58\%           &
20.13\%          & 0.239           & \textbf{62.84\%}  & \textbf{77.84\%} &
0.191           & 35.74\%           & 66.13\%          & 0.147           \\
\hline
SSDLite+MNetV2                       & \textbf{838ms} & 67.03\%           &
44.42\%          & 0.166           & 75.32\%           & 87.72\%          &
0.167           & 47.14\%           & 78.89\%          & 0.147           \\
Efficientdet-D0                      & 4686ms         & \textbf{73.31\%}  &
\textbf{58.86\%} & 0.148           & \textbf{79.14\%}  & \textbf{90.12\%} &
0.150           & \textbf{47.35\%}  & \textbf{79.06\%} & 0.143           \\
\hline
\end{tabular}}
\end{center}
\label{table:precision_tvn}
\end{table}

The poor classifier performance on the Caltech dataset can be due to several
factors, such as the low variability of backgrounds, the size of animals,
camouflage, quality of images, among others. However, our experiments were not
designed to identify these nuances intrinsic to this dataset. On the other
hand, this may be an interesting topic to be investigated in future work.

\textbf{Latency-precision trade-off:} Table~\ref{table:precision_tvn} shows
the latency for the inference of each model. Although EfficientDet-D0 offers a
strong baseline for the problem, its latency is more than four seconds, which is
prohibitive because camera trap images are usually obtained within one second
between them. In this context, SSDLite+MobileNetV2 may be deemed to attain
superior performance since its latency is below one second and it eliminated
at least 10\% more images than the classifiers with comparable latencies.

\subsection{Training with more Images}

A common way to improve model performance is to use more training
instances~\citep{bit2020Kolesnikov}. Following this strategy, an experiment was
carried out to train the classifiers using the subsets of the SS dataset. This
dataset contains ten times more images than the one used in the previous
experiment. Figure~\ref{fig:pr_curve_more_images} and
Table~\ref{table:more_images} summarize the results attained in this scenario.
The results indicate that classifiers outperformed detectors. We can highlight
that MobileNetV2 (224) reached performance similar to SSDLite+MobileNetV2, with
less than half of inference latency though. It is important to note, however,
that all detectors were trained using data subsets from the previous
experiment. Therefore, in order to provide a fair comparison, a similar amount
of training instances used to train the classifiers should be used to train the
detectors. On the other hand, obtaining more annotated instances for the
detection problem is expensive. This is the reason we do not show results using
the same dataset for training classifiers and detectors. However, we can
conclude based on the results obtained in this experiment that, depending on the
number of labeled instances available, classifiers can be a viable option for
the problem of identifying empty images, also showing better efficiency in terms
of computational resources.

\begin{table}[htb]
\caption[Comparison of precision for the nonempty class and the true negative
rate (TNR)]{Comparison of precision for the nonempty class and the true negative
rate (TNR), with a recall at 96\%, for classifiers trained in sets with ten
times more images from the SS dataset. Note the significant improvement in the
performance of the classifiers, exceeding by a large margin the detectors
trained in the original set.}
\begin{center}
\setlength{\tabcolsep}{3pt}
\begin{tabular}{l|cc|cc}
 \hline
\textbf{}                            & \multicolumn{2}{c|}{\textbf{SS-Site}}

                                & \multicolumn{2}{c}{\textbf{SS-Time}}
                          \\
\multicolumn{1}{c|}{\textbf{Model}} & \multicolumn{1}{c}{\textbf{Precision}} &
\multicolumn{1}{c|}{\textbf{TNR}} & \multicolumn{1}{c}{\textbf{Precision}} &
\multicolumn{1}{c}{\textbf{TNR}} \\ \hline \hline
Efficientnet-B0                      & 73.92\%                               &
86.78\%                           & 48.81\%                               &
80.25\%                          \\
Efficientnet-B3                      & \textbf{87.67\%}                      &
\textbf{94.73\%}                  & \textbf{64.28\%}                      &
\textbf{89.54\%}                 \\
MobileNetV2-224                      & 75.18\%                               &
87.63\%                           & 49.12\%                               &
80.50\%                          \\
MobileNetV2-320                      & \textbf{82.89\%}                      &
\textbf{92.26\%}                  & \textbf{58.61\%}                      &
\textbf{86.70\%}                 \\ \hline
\end{tabular}
\end{center}
\label{table:more_images}
\end{table}

\subsection{Model Optimization}

It is possible to observe from the results shown in previous sections that models
dealing with higher image input resolutions achieved higher performance but at
the expense of higher inference latency. In order to reduce latency of these
models, we performed an experiment using the SS dataset and MobileNetV2 (320).
In this experiment, the number of filters was reduced by adjusting the alpha
parameter. We also evaluated the post-training quantization of the models.
Since the swish activation function used by EfficientNets is not very well
supported for quantization~\citep{xiong2020mobiledets}, models of this family
were not used here. In addition, models evaluation was carried out on a sample
of 5,000 instances obtained from the validation set, because the TensorFlow Lite
quantized models are not optimized for inference on x86 architectures and
graphics cards used for training.

The results shown in Figure~\ref{fig:latency_prauc_ss} reinforce the role of
input resolution, since a MobileNetV2 with half the filters (MobileNetV2-0.50)
but higher resolution (320) obtained overall performance superior to the
performance of the standard model (MobileNetV2 (224)). In
Table~\ref{table:resolution}, we may observe that it is possible to reduce
about 23\% of model latency due to both quantization (SS-site) and number of
filters reduction (SS-time), maintaining compatible performance. It is worth
mentioning that the quantized version of the SSDLite + MobileNetV2 detector did
not obtain good results at the 96\% recall level used. These results indicate
the importance of assessing the performance of the models after quantization,
especially when the confidence threshold is adjusted at certain levels. One way
to avoid this problem would be quantization-aware training.

\begin{table}[htb]
\caption[Performance comparison of MobileNetV2 models for various widths on
the SS dataset]{Performance comparison of MobileNetV2 models for various widths
on the SS dataset. The models were converted to TensorFlow Lite and evaluated
using a sample of 5,000 instances randomly chosen from the validation set.}
\begin{center}
\resizebox{\textwidth}{!}{
\begin{tabular}{ll|c|c|cc|cc}
\hline
\textbf{}       &      &             &                &
\multicolumn{2}{c|}{\textbf{SS-Site}}                                  &
\multicolumn{2}{c}{\textbf{SS-Time}}                                 \\
\multicolumn{2}{c|}{\textbf{Model}} & \textbf{CPU} & \textbf{Memory}  &
\multicolumn{1}{c}{\textbf{Precision}} & \multicolumn{1}{c|}{\textbf{TNR}} &
\multicolumn{1}{c}{\textbf{Precision}} & \multicolumn{1}{c}{\textbf{TNR}} \\
\hline \hline
\multirow{3}{*}{MobileNetV2-224}  & Float       & 322ms  &  22.7MB    & 58.20\%




         & 73.07\%                           & 25.67\%

 & 45.33\%                          \\
                      & Int8        & \textbf{237ms} & \textbf{9.3MB} & 55.52\%


         & 69.99\%                           & 26.18\%

 & 46.74\%                          \\
                      & Int8 (quant aware)        & 239ms & 9.8MB
 & 57.31\% & 72.10\%  & - & -   \\ \hline
\multirow{2}{*}{MobileNetV2-0.50-320} & Float       & 259ms  & 19.7MB     &
54.86\%



         & 69.15\%                           & 27.76\%

 & 50.86\%                          \\
                      & Int8        & 237ms  &  9.6MB        & 51.50\%


         & 64.70\%                           & 26.67\%

 & 48.09\%                          \\ \hline
\multirow{3}{*}{MobileNetV2-0.75-320} & Float       & 484ms & 19.8MB   &
64.89\%


         & 79.72\%                           & \textbf{31.37\%}

 & \textbf{58.71\%}                 \\
                      & Int8        & 408ms   & 13.3MB       & 59.50\%


         & 74.49\%                           & 30.44\%

 & 56.87\%                          \\
                      & Int8 (quant aware)        & 413ms & 13.4MB
 & 61.38\% & 76.41\%  & - & -   \\ \hline
\multirow{2}{*}{MobileNetV2-320} & Float       & 635ms  & 33.5MB    & 66.65\%




         & 81.25\%                           & 30.34\%

 & 56.65\%                          \\
                      & Int8        & 485ms & 13.9MB  & \textbf{66.70\%}
         & \textbf{81.28\%}                  & 30.98\%

 & 57.95\%                          \\ \hline
\multirow{2}{*}{SSDLite+MobileNetV2} & Float       & 840ms & 31.7MB   & 73.68\%





         & 86.62\%                           & 39.50\%

 & 71.09\%                          \\
                      & Int8        & 575ms  &  13.6MB       & 36.83\%

         & 35.72\%                           & 19.27\%

 & 20.90\%                          \\ \hline
\end{tabular}}
\end{center}
\label{table:resolution}
\end{table}


To assess this, we fine tuned the models MobileNetV2-224 and
MobileNetV2-0.75-320 trained on SS-Site using quantization-aware training
for 2 epochs with the initial learning rate divided by 10. In this case, we were
able to improve the results of the quantized versions, as shown in
Table~\ref{table:resolution}. Unfortunately, we were not able to assess SSDLite
in the same way due to the fact that Tensorflow Object Detection API did not
support quantization-aware training for TF2 at the experiments' time.

Although the focus of this work is the latency-precision trade-off, we also
measured the memory usage of each model, as shown in
Table~\ref{table:resolution}. Quantized models have a significant reduction in
memory usage compared to the floating point models. However, when observed the
amount of RAM available on Raspberry Pi (1GB), the memory required by the
most complex model is very low (33.5MB), therefore, not affecting the
performance of other processes that may be running on the device.

\section{Final Remarks}

In this chapter, we presented a comparative study between detection and
classification models in the context of identification of nonempty images of
animals on edge devices. Our results showed that detection models outperform
classifiers when both are trained using the same training set, but their
superior inference latency can limit their use on edge devices. Moreover,
depending on the dataset, it is possible to train classifiers to obtain
satisfactory results, especially when there is massive number of images
available, as in the Snapshot Serengeti dataset. When detector is essential,
e.g. Caltech dataset, but the model's latency is not within the design
requirements, it may be necessary to use hardware accelerators, such as
EdgeTPUs or DSPs. Another limitation to detectors is the difficulty on obtaining
new instances annotated with bounding boxes, which is time-consuming and
expensive. In this situation, one possibility is to use techniques tackling few
or no labels, such as semi-supervised learning and self-supervised learning.
Other alternative is training an agnostic detection model designed to run on
edge devices, inspired by MegaDetector. Regarding the optimization strategies
evaluated, the post-training quantization proved to be effective, but it
requires a careful evaluation of the resulting model, which may decrease its
performance due to the quantization process. However, quantization-aware
training may solve this issue. Finally, using models with fewer filters but
with higher resolution was also effective. Therefore, techniques such as
knowledge distillation and model pruning are interesting directions for future
work.

In this chapter, we handled the filtering of empty images considering a two-stage
approach, where in the first stage a binary classifier distinguishes empty from 
nonempty images, leaving the species identification to the second stage. In the next 
chapter, the empty class is incorporated directly into the species classifier as 
just another category. As previously highlighted, empty images usually represent the
majority category in camera-trap datasets, with the remaining species categories
exhibiting a highly imbalanced, long-tail distribution. While we balanced the 
Snapshot Serengeti dataset in this chapter to better handle the filtering task, in 
the next chapter, we keep the original dataset distributions. In fact, the next 
chapter treats this extreme imbalance as an inherent characteristic of camera-trap 
data, focusing on improving performance under this constraint.
